% ID: SEL-D0075-OB-E03
% ESTADO: cotejada. Segunda pasada visual y borrador compilado revisados.
% Fuente: Selectividad/sel_2008_matematicas_aplicadas/Mat CCSS 3 - Examen y criterios de correción.pdf
% SHA-256: ebf609a5063d63be15b900270973808e3fe0a239a4d911d6b41d80544711aeae
% Página: 2; opción B; ejercicio original 3.

\noindent\underline{Parte I}

Se sabe que el $30\%$ de los individuos de una población tiene estudios superiores; también se sabe que, de ellos, el $95\%$ tiene empleo. Además, de la parte de la población que no tiene estudios superiores, el $60\%$ tiene empleo.
\begin{enumerate}[label=\alph*)]
\item \textbf{(1 punto)} Calcule la probabilidad de que un individuo, elegido al azar, tenga empleo.
\item \textbf{(1 punto)} Se ha elegido un individuo aleatoriamente y tiene empleo; calcule la probabilidad de que tenga estudios superiores.
\end{enumerate}
\par\noindent\begin{minipage}{\linewidth}
\noindent\underline{Parte II}

Sea la población $\{1,2,3,4\}$.
\begin{enumerate}[label=\alph*)]
\item \textbf{(1 punto)} Construya todas las muestras posibles de tamaño 2, mediante muestreo aleatorio simple.
\item \textbf{(1 punto)} Calcule la varianza de las medias muestrales.
\end{enumerate}
\end{minipage}\par
